\documentclass[10pt,a4paper]{amsart}
\usepackage[margin=19mm]{geometry}
\usepackage{amsmath,amssymb,mathtools}
\usepackage[scr=rsfs,cal=euler]{mathalfa}
\usepackage{xfrac}
\usepackage[unicode,hidelinks,bookmarks=false]{hyperref}
\newcommand{\ie}{i.e.\ }
\newcommand{\cf}{cf.\ }
\newcommand{\resp}{respectively\ }
\newcommand{\Subsection}{\subsection}
\providecommand{\nobreakdash}{\nobreak\hbox{-}}
\setlength{\emergencystretch}{2em}
\multlinegap0pt
\usepackage{bm}
\usepackage{bbm}
\usepackage{dsfont}
\usepackage{xspace}
\usepackage{cancel}
\usepackage{mathtools}
\usepackage{amsthm}
\makeatletter
\@ifundefined{theorem}{\newtheorem{theorem}{Theorem}}{}
\makeatother
\def\AND{\, \wedge \,}
\def\add{\text{add}}
\def\suchthat{\ : \ }
\title[Complexity of Linear Subsequences of $k$-Automatic Sequences]{Complexity of Linear Subsequences of $k$-Automatic Sequences}
\author{Delaram Moradi, Narad Rampersad,, Jeffrey Shallit}
\date{Source: arXiv:2512.10017v5 (CC BY 4.0)}
\begin{document}
\maketitle

\noindent\emph{Source and licence.} Complexity of Linear Subsequences of $k$-Automatic Sequences. Version arXiv:2512.10017v5 (CC BY 4.0), distributed under the Creative Commons Attribution 4.0 International licence, as recorded in the arXiv licence metadata for this exact version and shown on the abstract page https://arxiv.org/abs/2512.10017v5. Full rights notice: licenses/arxiv-2512.10017-CC-BY-4.0.txt.\par\smallskip

\noindent\textit{Editorial scope.} Counted unit: the Theorem whose source label is \texttt{theorem:y=n[x]-time} in the article (source0.tex lines 1095--1098, source label \texttt{theorem:y=n[x]-time}), delivered together with the article's own complete original proof of it (source0.tex lines 1100--1164, one \texttt{proof} environment of 65 lines). No mathematics is written, repaired or re-derived: statement and proof are the article's own text line for line. Required context retained uncounted: theorem:n[x]+c=[y]-sc (lines 391--394). Every definition, display and label of those lines is retained unchanged. Omitted from this package and not redistributed: 1--390, 395--1094, 1099--1099, 1165--1281. Nothing inside a retained excerpt is omitted or altered: every retained body is delivered exactly as the article's lines are written, so no omission receipt of any kind accompanies this package. Citations: the retained excerpts carry no citation command at all, so no bibliography entry of the article is transcribed and no reference is redistributed. Reference adaptation: none was needed and none was made; every reference inside the delivered unit resolves inside it, so no unresolved reference or citation is produced. Printed slips kept literal: no printed word, symbol, hypothesis or number of the counted statement or of its proof was altered. Layout and class: the article's own class is \texttt{article} with packages including amssymb, amsmath, amsthm, amsfonts, amscd, graphicx, color, hyperref, algorithm2e, xcolor, mdframed, enumitem, fullpage, float; the wrapper is the lane's pinned \texttt{amsart} editorial wrapper at 10pt on A4, which restates the theorem-like environments the retained text needs and adds the article's own macro definitions that the retained text needs (\texttt{\textbackslash AND}, \texttt{\textbackslash add}, \texttt{\textbackslash suchthat}), with extra wrapper packages (bm, bbm, dsfont, xspace, cancel, mathtools). The delivered numbering is the unit's own. Assets: no retained span contains a figure environment, a graphics-inclusion command, a PDF-page inclusion, a table or a TikZ picture; no image file is copied into this package, the package declares no asset dependency and no image file is redistributed with it. Licence: version arXiv:2512.10017v5 of this article is distributed under the Creative Commons Attribution 4.0 International licence, as recorded in the arXiv licence metadata for this exact version and shown on the abstract page https://arxiv.org/abs/2512.10017v5. A full rights notice is delivered at licenses/arxiv-2512.10017-CC-BY-4.0.txt. Characters: every retained excerpt is pure ASCII, so the delivered engine, XeLaTeX with its default Latin Modern text font, reports no missing character.
\begin{theorem} \label{theorem:n[x]+c=[y]-sc}
Let $n \geq 1$, $0 \leq c < n$ be fixed constants.
For both lsd- and msd-first input, there is an $n$-state automaton accepting $x \times y$ where $n[x]_k + c = [y]_k$.
\end{theorem}

\begin{theorem} \label{theorem:y=n[x]-time}
    Let $n \geq 1$ be a fixed constant.
    The automaton for recognizing $[y]_k = n[x]_k$ on msd-first input can be computed in $O(n \log^2 n)$ time. 
\end{theorem}

\begin{proof}
First by induction on $n$ we prove that the we can recursively create DFA $M_n$ with msd-first input recognizing $[y]_k = n[x]_k$ with at most $n$ states.

Base case: $n=1$.

In this case, a simple $1$-state automaton recognizes $[y]_k=[x]_k$. The single state corresponds to $[y]_k - [x]_k = 0$.

For the induction step, assume the result is true for all $n'<n$; we prove it for $n$.
There are two cases.

\bigskip\noindent{\it
Case 1:} $n$ is even.

By the induction hypothesis, we have a DFA $M_{n/2}$ with $n/2$ states corresponding to differences $[y]_k - (n/2)[x]_k$ in the range $[0, (n/2)-1]$.
In this case, we use the formula
$$ \exists y \ M_{n/2}(x,y) \AND M_{\add}(y, y, z).$$
First, we create an automaton $M'_{n}$ that is the product construction of $M_{n/2}$ and $M_{\add}$. So this automaton has $2(n/2) = n$ states. Consider an input $x \times y \times z$ for $M'_{n}$ leading to some state $[q, p]$ where $q$ is from $M_{n/2}$ and $p$ is from $M_{\add}$. By the construction of $M'_{n}$, we know $q=[y]_k - (n/2)[x]_k$ and $p=[z]_k - ([y]_k + [y]_k)$. So we have $2q+p = [z]_k - n[x]_k$. Since $p \in \{0, 1\}$ and $q \in \{0, \ldots, n/2 - 1\}$, all differences in $\{0, \ldots, n-1\}$ are present in $M'_{n}$ and they are only represented by a single state.
Therefore, the automaton $M'_{n}$ keeps track of the difference $[z]_k - n[x]_k$ on input $x \times y \times z$. The difference $[z]_k - n[x]_k$ that the automaton keeps track of is independent of $y$.

Then we apply the $\exists y$ quantifier by removing the component corresponding to $y$ from the transitions and we get an NFA. In the NFA there are no two same transitions from a single state leading to different states. So after subset construction the resulting DFA is the same as the NFA and we get the $n$-state automaton $M_{n}=(Q, \Sigma^2_k, \delta, q_0, A)$ as follows.
\begin{align*}
Q &= \{ d \in \mathbb{N}  \suchthat d \in [0, n-1] \}, \\
\delta(d, [a, e]) &= kd + e - n
a, \\
q_0 &= 0, \\
A &= \{q_0\}.
\end{align*}
Projecting away the $y$ component from transitions could have caused nondeterminism, but this is not the case.

\bigskip\noindent{\it
Case 2:} $n$ is odd.

By the induction hypothesis, we have a DFA $M_{n-1}$ with $n-1$ states corresponding to $[y]_k - (n-1)[x]_k$ in the range $[0, n-2]$.
In this case, we use the formula
$$ \exists y \ M_{n-1}(x,y) \AND M_{\add}(x, y, z).$$
First, we create an automaton $M'_{n}$ from the usual product construction for $M_{n-1}$ and $M_{\add}$. So this automaton has $2(n-1) $ states. Consider an input $x \times y \times z$ for $M'_{n}$ leading to some state $[q, p]$ where $q$ is from $M_{n-1}$ and $p$ is from $M_{\add}$. By the construction of $M'_{n}$, we know $q=[y]_k - (n-1)[x]_k$ and $p=[z]_k - ([x]_k + [y]_k)$. So we have $q+p = [z]_k - n[x]_k$. Therefore, to each state in the automaton $M'_{n}$ reachable on $x \times y \times z$, a difference $d=[z]_k - n[x]_k$ can be attributed. Furthermore, since $p \in \{0, 1\}$ and $q \in \{0, \ldots, n-2\}$, all differences in $\{0, \ldots, n-1\}$ are present in $M'_{n}$. Note that the difference $0$ corresponds to exactly one pair---namely $[0,0]$---and the difference $n-1$ corresponds to only one pair---namely $[n-2,1]$.
All the other differences $d$ (those in $\{ 1,2,\ldots, n-2\}$)  are represented by two states---namely, $[d-1,1]$ and $[d,0]$.

Next, we project away the $y$ component.  The resulting automaton is an NFA, and we use the subset construction to create a DFA.  Consider a state from $M_{n}$ reachable on $x \times z$. This state consists of some pairs $[q, p]$ from the NFA such that $[z]_k - n[x]_k = q+p$. As mentioned above, if $[z]_k-n[x]_k$ is $0$ or $n-1$, then there is only one such pair; otherwise, there are two such pairs.
 So the DFA has $O(n)$ states and the determinization did not cause a blow up in the number of states.

From the proof of Theorem \ref{theorem:n[x]+c=[y]-sc}, we know the automaton recognizing $n[x]_k= [z]_k$ on input $x \times z$ only needs to keep track of a difference $[z]_k - n[x]_k$ in range $[0, n-1]$. So after some minimization, the DFA  $M_{n}=(Q, \Sigma^2_k, \delta, q_0, A)$ can be formally defined as follows.
\begin{align*}
Q &= \{ d \in \mathbb{N}  \suchthat d \in [0, n-1] \}, \\
\delta(d, [a, e]) &= kd + e - na,
 \\
q_0 &= 0, \\
A &= \{q_0\}.
\end{align*}

Let $T(n)$ be the time used by the algorithm above to create $M_n$. We now analyze $T(n)$ based on the information above.

In Case 1, the number of transitions traversed for creating $M_n$ is as follows: $O(1)$ transitions for $M_{\add}$, $O(n)$ transitions for the product construction of $M_{n/2}$ and $M_{\add}$, $O(n)$ transitions for projecting away the $y$ component, $O(n)$ transitions for the determinized automaton, giving a total of $O(n)$ transitions in addition to the number of transitions traversed to create $M_{n/2}$.

In Case 2, the number of transitions traversed for creating $M_n$ is as follows: $O(1)$ transitions for $M_{\add}$, $O(n)$ transitions for the product construction of $M_{n-1}$ and $M_{\add}$, $O(n)$ transitions for projecting away the $y$ component, $O(n)$ transitions for the determinized automaton, giving a total of $O(n)$ transitions in addition to the number of transitions to create $M_{n-1}$.

Putting everything together and taking minimization steps into account, we get the following recursive formula:
\[ T(n) = 
\begin{cases}
    O(n \log n) + T(n/2), & \text{if $n \bmod 2 = 0$}; \\
    O(n \log n) + T(n-1), & \text{if $n \bmod 2 = 1$}.
\end{cases}
\]
So we have $T(n) = O(n\log^2 n)$.    
\end{proof}

\end{document}
